\section{Заключение}

В ходе выполнения выпускной квалификационной работы была достигнута поставленная цель: разработана расширяемая система генерации и диспетчеризации событий, ориентированная на интеграцию в системы глубокого анализа сетевого трафика (DPI). Для решения поставленных задач был проведён сравнительный анализ существующих анализаторов трафика и библиотек обработки событий, выявлены их ограничения в части масштабируемости, детерминированности задержек и гибкости расширения.

На основе полученных результатов спроектирована и реализована архитектура системы, обеспечивающая типобезопасное расширение через наследование от параметризованного базового класса \texttt{EventBase} с использованием паттерна CRTP. Ключевой особенностью разработанного решения стал механизм идентификации типов событий на этапе выполнения без применения RTTI или виртуальных функций, что позволило достичь нулевых накладных расходов на получение идентификатора в горячем пути диспетчеризации. Для хранения разнотипных обработчиков реализован механизм type erasure на базе библиотеки \texttt{function2} с оптимизацией small buffer optimization (SBO), исключающий динамическое выделение памяти для подавляющего большинства сценариев использования.

Разработана гибридная модель диспетчеризации, явно разделяющая синхронное выполнение CPU-bound обработчиков и асинхронную обработку IO-bound задач через выделенный runtime. Интерфейс системы обеспечивает проверку корректности на этапе компиляции за счёт использования концептов C++20 и идеальной передачи аргументов, что минимизирует риск ошибок при интеграции пользовательских модулей анализа.

Проведённые экспериментальные исследования подтвердили эффективность предложенных архитектурных и алгоритмических решений.
Базовая стоимость диспетчеризации события без обработчиков составила 2,1--2,4 нс, а с одним синхронным обработчиком --- 4,7 нс.
Система продемонстрировала почти линейное масштабирование при увеличении числа потоков обработки (эффективность 95\% при 4 потоках и 90\% при 8 потоках).
Сравнительный анализ показал преимущество разработанной системы над альтернативными решениями (Boost.Signals2, Eventpp, Sigslot) по задержке выполнения (в 4--12 раз и пропускной способности.

Практическая значимость работы заключается в создании готового к использованию программного компонента, который может быть интегрирован в существующие DPI-платформы, средства мониторинга сетевой безопасности и системы обработки данных в реальном времени.
Разработанная система событий представляет собой эффективное, масштабируемое и типобезопасное решение, удовлетворяющее требованиям современных систем глубокого анализа сетевого трафика и готовое к применению в производственной эксплуатации.
\newpage